\chapter{A polynomial separation between adaptive and non-adaptive algorithms}
\label{chap:separation}

This chapter formalizes Section~2.5 and Appendix~I of the paper: an action set
$\Xs \subset \R^{kd}$, a union of $k$ unit balls living in $k$ disjoint blocks of $d$
coordinates, on which every non-adaptive algorithm needs $\Omega(kd^2/\eps^2)$ samples while
an adaptive algorithm succeeds with $O(kd\log(k/\delta)/\eps^2 + d^2/\eps^2)$ samples. The
target is Theorem~7 of the paper (Theorem~\ref{thm:separation}). The adaptive algorithm
first runs the norm estimation procedure of Chapter~\ref{chap:norm} on each block to locate
the block containing a near-optimal arm, then identifies a good arm in that block with the
basis design of Lemma~\ref{lem:ball_identification}; as in Chapter~\ref{chap:norm}, the
phases are composed by a deterministic round schedule of a seeded algorithm
(Lemma~\ref{lem:seeded}), and the analysis takes place on the seed space. The chapter is
fully formalized. The non-adaptive lower bound deviates from the paper: instead of
the cited adaptive lower bound of Chen et al.\ \cite{chen2024assouad} we project a fixed
design on the block receiving the fewest points (Lemma~\ref{lem:block_restriction}, an
instance of the transport Lemma~\ref{lem:fixed_design_transport}) and apply the ingredients
of our own non-adaptive lower bound (the Bayesian step,
Lemma~\ref{lem:lower_nonadaptive_bayes_step}) with the width of the unit ball for a matrix
(Proposition~\ref{prop:width_ball}, $\gwidth{\ball_d}{A} \ge \sqrt{2/\pi}\,d/\sqrt{\tr A}$). The class of
non-adaptive algorithms is therefore that of Definition~\ref{def:fixed_design}: deterministic
fixed designs with an arbitrary (possibly randomized) recommendation kernel, whereas the
paper's Theorem~7 speaks of ``possibly randomized non-adaptive'' algorithms
(Remark~\ref{rem:separation_gap}). The block-ball set is spanning
(Lemma~\ref{lem:block_ball_basic}), which is the hypothesis needed wherever
Theorem~\ref{thm:lower_nonadaptive}, Theorem~\ref{thm:lower_adaptive} or the Gaussian width
(Definition~\ref{def:width}, Proposition~\ref{prop:width_ball} on the spanning set $\ball_d$)
is used below. Throughout,
$k, d \ge 1$ are integers with $kd \ge 2$; the comparison of the two bounds uses
$d \le k \le d^2$.

\section{The block-ball action set}
\label{sec:separation_set}

\begin{definition}[Block-ball action set]\label{def:block_ball_set}
\lean{Maiti2026Power.blockBallSet}\leanok
\uses{def:arm_set}
For $i \in \{1,\dots,k\}$ let $B_i := \{(i-1)d + 1, \dots, id\}$ be the $i$-th block of
coordinates of $\R^{kd}$, let $\iota_i : \R^d \to \R^{kd}$ be the isometric embedding which
places a vector on the coordinates of $B_i$ and $0$ elsewhere, and let
$\pi_i : \R^{kd} \to \R^d$ be the projection on the coordinates of $B_i$, so that
$\pi_i\circ\iota_i = \mathrm{id}$, $\ip{\iota_i u}{v} = \ip{u}{\pi_i v}$ and
$\norm{\iota_i u} = \norm u$. Define
\[
  \Xs_i := \{ x \in \R^{kd} : \supp(x) \subseteq B_i,\ \norm{x} \le 1 \} = \iota_i(\ball_d),
  \qquad
  \Xs := \bigcup_{i=1}^k \Xs_i .
\]
For $\theta \in \R^{kd}$ we write $\theta^{(i)} := \pi_i\theta \in \R^d$ for its $i$-th block,
so $\theta = \sum_i\iota_i\theta^{(i)}$ and $\norm{\theta}^2 = \sum_i\norm{\theta^{(i)}}^2$.
\end{definition}

\textbf{Lean remark.} With $\R^{kd}$ encoded as \texttt{EuclideanSpace ℝ (Fin k × Fin d)}
the blocks are $\{i\}\times\mathrm{Fin}\ d$ (0-indexed) and $\iota_i$, $\pi_i$ are the
obvious linear maps; $\Xs_i \cap \Xs_j = \{0\}$ for $i \ne j$, so an element of $\Xs$
determines its block unless it is $0$. Where a block index is needed for a point of $\Xs$
we use the smallest one.

\begin{lemma}[Basic properties of the block-ball set]\label{lem:block_ball_basic}
\lean{Maiti2026Power.blockEmb, Maiti2026Power.blockProj, Maiti2026Power.inner_blockEmb_left, Maiti2026Power.norm_blockEmb, Maiti2026Power.norm_blockProj_le, Maiti2026Power.blockProj_blockEmb_same, Maiti2026Power.blockProj_blockEmb_of_ne, Maiti2026Power.exists_blockProj_eq_zero, Maiti2026Power.blockBallSet_eq_iUnion, Maiti2026Power.isCompact_blockBallSet, Maiti2026Power.blockBallSet_subset_unitBall, Maiti2026Power.span_blockBallSet, Maiti2026Power.supportFn_blockBallSet, Maiti2026Power.supportFn_blockBallSet_blockEmb, Maiti2026Power.simpleRegret_blockBallSet_blockEmb, Maiti2026Power.supportFn_unitBall, Maiti2026Power.simpleRegret_unitBall}\leanok
\uses{def:block_ball_set, def:regret}
$\Xs = \bigcup_i\iota_i(\ball_d)$ is a compact subset of $\ball_{kd}$ with $\Span(\Xs) = \R^{kd}$,
hence an action set in the sense of Definition~\ref{def:arm_set}. Every $x \in \Xs$ satisfies
$\pi_j x = 0$ for all $j$ but one. For every $\theta \in \R^{kd}$,
\[
  \max_{x \in \Xs}\ip{x}{\theta} = \max_{1\le i\le k}\norm{\theta^{(i)}},
\]
and for $\vartheta \in \R^d$ and $\rec \in \R^{kd}$,
$\max_{x\in\Xs}\ip{x}{\iota_i\vartheta} = \norm{\vartheta} = \max_{u \in \ball_d}\ip{u}{\vartheta}$
and
\[
  r_{\Xs}(\rec, \iota_i\vartheta) = \norm{\vartheta} - \ip{\pi_i\rec}{\vartheta}
  = r_{\ball_d}(\pi_i\rec, \vartheta) .
\]
\end{lemma}

\begin{proof}
\leanok
\uses{def:block_ball_set}
Each $\iota_i(\ball_d)$ is the continuous image of a compact set, and a finite union
of compact sets is compact; $x \in \Xs$ lies in $\iota_i(\ball_d)$ for its block $i$, with
$x = \iota_i\pi_i x$. $\Xs$ contains every basis vector $e_m$ ($m \in B_i$ for some
$i$), so it spans $\R^{kd}$; and $\norm{\iota_i u} = \norm u \le 1$ gives $\Xs \subseteq \ball_{kd}$.
For $x = \iota_i u$ with $u \in \ball_d$, $\ip{x}{\theta} = \ip{u}{\theta^{(i)}} \le \norm{\theta^{(i)}}$
by Cauchy--Schwarz, with equality for $u = \theta^{(i)}/\norm{\theta^{(i)}}$ (any $u$ if
$\theta^{(i)} = 0$). Taking the maximum over $i$ gives the second identity; for
$\theta = \iota_i\vartheta$ the blocks are $\theta^{(i)} = \vartheta$ and $\theta^{(j)} = 0$
for $j \ne i$, and the same Cauchy--Schwarz argument on $\ball_d$ gives
$\max_{u\in\ball_d}\ip{u}{\vartheta} = \norm\vartheta$. The regret formula follows from
Definition~\ref{def:regret} and $\ip{\rec}{\iota_i\vartheta} = \ip{\pi_i\rec}{\vartheta}$.
\end{proof}

\section{Lower bound for non-adaptive algorithms}
\label{sec:separation_lower}

\begin{lemma}[Restriction of a fixed design to a block]\label{lem:block_restriction}
\lean{Maiti2026Power.blockProjBall, Maiti2026Power.blockDesign, Maiti2026Power.blockRounds, Maiti2026Power.isPAC_fixedDesignTransport_blockProjBall, Maiti2026Power.trace_sum_outerSelf_blockDesign_le, Maiti2026Power.pairwiseDisjoint_blockRounds, Maiti2026Power.sum_card_blockRounds_le, Maiti2026Power.exists_mul_card_blockRounds_le}\leanok
\uses{def:fixed_design, def:block_ball_set, lem:block_ball_basic, def:pac, lem:fixed_design_transport}
Let $x_0,\dots,x_{T-1} \in \Xs$ be a fixed design with recommendation kernel $\rho$ (from
$\R^T$ to $\Xs$) which is $(\eps,\delta)$-PAC on $\Xs$, and fix a block $i$. The projected
design $x'_t := \pi_i x_t \in \ball_d$ ($t < T$), with the recommendation kernel
$(\pi_i)_*\rho$ (draw $\rec \sim \rho(y)$, the history being relabelled with the actions $x_t$,
and output $\pi_i\rec \in \ball_d$), is a fixed-design algorithm with budget $T$ on $\ball_d$
which is $(\eps,\delta)$-PAC on $\ball_d$. Its design matrix
$\Sigma'_T := \sum_{t<T}(\pi_ix_t)(\pi_ix_t)^\top$ satisfies $\tr\Sigma'_T \le T_i$, where
$T_i := \abs{\{t < T : \pi_i x_t \ne 0\}}$ is the number of rounds played in block $i$; the sets
$\{t < T : \pi_ix_t \ne 0\}$ are pairwise disjoint, so $\sum_i T_i \le T$ and some block $i^*$
has $kT_{i^*} \le T$.
\end{lemma}

\begin{proof}
\leanok
\uses{lem:fixed_design_transport, lem:block_ball_basic}
This is Lemma~\ref{lem:fixed_design_transport} with $L = \iota_i$ and $\pi = \pi_i$: by
Lemma~\ref{lem:block_ball_basic}, $\pi_ix_t \in \ball_d$ (as $\norm{\pi_i x} \le \norm x \le 1$),
$\ip{\pi_ix_t}{\vartheta} = \ip{x_t}{\iota_i\vartheta}$ (adjointness) and
$r_{\ball_d}(\pi_i\rec,\vartheta) = r_{\Xs}(\rec,\iota_i\vartheta)$ for every $\vartheta \in \R^d$.
The trace bound is $\tr\Sigma'_T = \sum_{t<T}\norm{\pi_ix_t}^2 \le T_i$ since each term is at
most $1$ and vanishes outside the rounds played in block $i$. The sets of rounds are disjoint
because each $x_t$ lies in a single block (Lemma~\ref{lem:block_ball_basic}), so
$\sum_i T_i \le T$, and the pigeonhole principle gives $i^*$ with $kT_{i^*} \le T$.
\end{proof}

\textbf{Lean remark.} The paper (and an earlier version of this blueprint) restricts the
design to the rounds $t$ with $x_t$ in block $i$; keeping all $T$ rounds with the projected
actions $\pi_ix_t$ (which are $0$ outside the block and contribute pure noise) gives a
fixed-design algorithm with the same budget $T$, which is exactly the situation of
Lemma~\ref{lem:fixed_design_transport}, and the number $T_i$ of rounds in the block enters
only through the trace of the design matrix. The composition of $\rho$ with the relabelling of
the history and with $\pi_i$ is where the randomness allowed in Definition~\ref{def:bai_alg}
is used.

\begin{theorem}[Non-adaptive lower bound for the block-ball set]\label{thm:separation_lower_nonadaptive}
\lean{Maiti2026Power.blockBallSet_le_budget_of_isFixedDesign_of_isPAC, Maiti2026Power.blockBallSet_le_budget_of_isPAC}\leanok
\uses{def:block_ball_set, lem:block_ball_basic, def:fixed_design, def:pac}
Let $\eps \in (0,1]$. Recall that $\Xs$ is a spanning action set in $\R^{kd}$
(Lemma~\ref{lem:block_ball_basic}).
\begin{enumerate}
  \item[(i)] If a fixed-design algorithm (Definition~\ref{def:fixed_design}: a deterministic
        design $x_0,\dots,x_{T-1} \in \Xs$ with an arbitrary, possibly randomized,
        recommendation kernel; budget $T \ge 0$) is $(\eps,\delta)$-PAC on $\Xs$ with
        $\delta \le 1/10$, then
        \[
          T \ge \frac{k d^2}{210\,\eps^2} .
        \]
  \item[(ii)] If $kd \ge 2$ and any identification algorithm (adaptive or not) with budget $T$
        is $(\eps,\delta)$-PAC on $\Xs$ with $\delta < 1/16$, then
        $T \ge kd\log(1/\delta)/(20000\,\eps^2)$ (and $T \ge c_{\mathrm{ad}}\,kd\log(1/\delta)/\eps^2$
        with the constant $c_{\mathrm{ad}}$ of Theorem~\ref{thm:lower_adaptive}).
\end{enumerate}
\end{theorem}
In particular every fixed-design $(\eps,\delta)$-PAC algorithm with $\delta < 1/16$ has
budget at least
$\tfrac12\big( c_{\mathrm{ad}}\,kd\log(1/\delta) + kd^2/210 \big)/\eps^2$.

\begin{proof}
\leanok
\uses{lem:block_restriction, lem:lower_nonadaptive_posdef, lem:lower_nonadaptive_bayes_step, prop:width_ball, thm:lower_adaptive, lem:block_ball_basic}
(i) If $k = 0$ or $d = 0$ there is nothing to prove, so assume $k, d \ge 1$. With the notation
of Lemma~\ref{lem:block_restriction}, pick a block $i^*$ with $kT_{i^*} \le T$ and consider
the projected design $x'_t = \pi_{i^*}x_t$ on $\ball_d$, with design matrix $\Sigma'_T$, which
is $(\eps,\delta)$-PAC on $\ball_d$ with budget $T$.

\emph{Step 1: $\Sigma'_T \succ 0$.} Since $0 \in \ball_d$, $\ball_d$ is spanning and
$\delta \le 1/10 < 1/2$, Lemma~\ref{lem:lower_nonadaptive_posdef} applies to the projected
design (with any budget $T \ge 0$, including $T = 0$). In particular $T_{i^*} \ge 1$ (as
$d \ge 1$ and $\tr\Sigma'_T \le T_{i^*}$, with $\tr\Sigma'_T > 0$).

\emph{Step 2: the Bayesian step.} Lemma~\ref{lem:lower_nonadaptive_bayes_step} on $\ball_d$
gives $0.12\,\gwidth{\ball_d}{\Sigma'_T} \le \eps$.

\emph{Step 3: the width of the ball.} By Proposition~\ref{prop:width_ball},
$\gwidth{\ball_d}{\Sigma'_T} \ge \sqrt{2/\pi}\,d/\sqrt{\tr\Sigma'_T} \ge \sqrt{2/\pi}\,d/\sqrt{T_{i^*}}$.

\emph{Step 4: arithmetic.} Combining, $0.12\sqrt{2/\pi}\,d \le \eps\sqrt{T_{i^*}}$; squaring,
$0.0144\cdot(2/\pi)\,d^2 \le \eps^2T_{i^*}$, and since $\pi \le 4$, $2/\pi \ge 1/2$ and
$0.0072\,d^2 \le \eps^2 T_{i^*}$, i.e.\ $T_{i^*} \ge d^2/(139\,\eps^2) \ge d^2/(210\,\eps^2)$.
Hence $T \ge kT_{i^*} \ge kd^2/(210\eps^2)$.

(ii) $\Xs$ is a spanning action set in $\R^{kd}$ with $kd \ge 2$
(Lemma~\ref{lem:block_ball_basic}), so Theorem~\ref{thm:lower_adaptive} applies directly with
dimension $kd$.

The last claim is $T \ge \max(a,b) \ge (a+b)/2$.
\end{proof}

\section{The adaptive block algorithm}
\label{sec:separation_algorithm}

\begin{lemma}[Transporting an algorithm on $\ball_d$ to a block]\label{lem:separation_block_embedding}
\lean{Maiti2026Power.blockEmbBall, Maiti2026Power.blockProjBall_blockEmbBall, Maiti2026Power.BlockParam.F_blockEmbBall, Maiti2026Power.BlockParam.seedStep_block, Maiti2026Power.BlockParam.yω_block, Maiti2026Power.BlockParam.estN_eq, Maiti2026Power.shiftSeq, Maiti2026Power.nsMeasure_map_shiftSeq}\leanok
\uses{def:block_ball_set, def:env, def:bai_alg, def:norm_estimator, lem:seeded, def:meta_algorithm}
Let $\iota_i : \ball_d \to \Xs$ be the embedding of the unit ball on block $i$ (so
$\pi_i \circ \iota_i = \mathrm{id}$). For every $\theta \in \R^{kd}$ and $x \in \ball_d$, the
reward kernel of $\nu_\theta$ at $\iota_i x$ is that of $\nu_{\theta^{(i)}}$ at $x$:
$\ip{\iota_i x}{\theta} + \eta = \ip{x}{\theta^{(i)}} + \eta$. Consequently, the block
algorithm of Definition~\ref{def:block_algorithm} \emph{simulates} the norm estimation
meta-algorithm on each block: on the seed space, for $i < k$ and $t < N$, the step
(action, observation) of round $iN + t$ of the block algorithm against $\nu_\theta$ is
$(\iota_i x_t, y_t)$ where $(x_t, y_t)$ is the step of round $t$ of
$\mathrm{NormEst}(\eps/4, \delta/(2k))$ against $\nu_{\theta^{(i)}}$ on the seed sequence shifted
by $iN$, $(\omega_{iN + t})_t$. In particular the estimate $\hat r_i$ of block $i$ is the
estimate of the meta-algorithm against $\nu_{\theta^{(i)}}$ on the shifted seed sequence, and the
shifted seed sequence has the law $\Prob_{\mathrm{seed}}$ of the seed sequence.
\end{lemma}

\begin{proof}
\leanok
\uses{def:env, lem:block_ball_basic, lem:seeded, lem:norm_meta_plan, lem:statistic_structure}
The kernel identity is Lemma~\ref{lem:block_ball_basic} ($\ip{\iota_i x}{\theta} = \ip{x}{\pi_i\theta}$).
The simulation is a strong induction on $t < N$: the action of round $iN + t$ of the block
algorithm is $\iota_i$ applied to the next action of the meta-algorithm computed from the past
actions (projected by $\pi_i$) and observations of the rounds $iN + r$, $r \le t$, of the
block, which by the induction hypothesis are the past actions and observations of the
meta-algorithm on the shifted seed sequence; the observation follows from the kernel
identity. The shift invariance of $\Prob_{\mathrm{seed}}$ is Lemma~\ref{lem:statistic_structure}
(the coordinates $(\omega_{a + t})_t$ of an i.i.d.\ sequence form an i.i.d.\ sequence with the
same law).
\end{proof}

\textbf{Lean remark.} The block algorithm is not obtained by transporting the meta-algorithm
as an LML algorithm (which would require mapping its policy kernels by $\iota_i$ and its
histories by $\pi_i$): it is a seeded algorithm whose next-action map, in the rounds of block
$i$, applies $\iota_i$ to the next-action map of the meta-algorithm evaluated on the
$\pi_i$-projected history of the block. The simulation lemma is then a statement about two
recursively defined runs on the same seed space.

\begin{proof}
\uses{def:env, lem:block_ball_basic}
The reward kernel of $\nu_\theta$ at $\iota_i(x)$ is
$\Normal(\ip{\iota_ix}{\theta}, 1) = \Normal(\ip{x}{\theta^{(i)}}, 1)$
(Lemma~\ref{lem:block_ball_basic}), which is the reward kernel of $\nu_{\theta^{(i)}}$ at $x$.
The policy of $\iota_i\mathrm{alg}$ is the image by $\iota_i$ of the policy of $\mathrm{alg}$
applied to the projected history, and $\pi_i\circ\iota_i = \mathrm{id}$, so the projected
process satisfies the defining conditional-law conditions of an algorithm-environment
sequence for $(\mathrm{alg},\nu_{\theta^{(i)}})$ (LML \texttt{IsAlgEnvSeq}); uniqueness of the
trajectory law gives the claim about estimates.
\end{proof}

\textbf{Lean remark.} This is the analogue, for the non-invertible isometric embedding
$\iota_i$, of the transformation lemma \ref{lem:transform_preserves_pac} used for the
adaptive lower bound; it is a statement about mapping actions by a measurable map $\phi$
such that the environments correspond ($\kappa_\theta \circ \phi = \kappa_{\theta'}$), which
holds for any stationary environments and could be stated at that level of generality. The
map $\phi : \ball_d \to \Xs$, $u \mapsto \iota_i u$, is continuous hence measurable between
the subtypes, and the policy kernels of $\iota_i\mathrm{alg}$ are the images of those of
$\mathrm{alg}$ by $\phi$ (\texttt{Kernel.map}), applied to the history mapped by $\pi_i$ on
its action coordinates; the action type of $\iota_i\mathrm{alg}$ is the subtype $\Xs$, as
required to compose it with other algorithms on $\Xs$ (Lemma~\ref{lem:composition}).

\begin{definition}[Block algorithm]\label{def:block_algorithm}
\lean{Maiti2026Power.BlockParam, Maiti2026Power.BlockParam.P, Maiti2026Power.BlockParam.N, Maiti2026Power.BlockParam.n₂, Maiti2026Power.BlockParam.T₂, Maiti2026Power.BlockParam.T, Maiti2026Power.BlockParam.blockOf, Maiti2026Power.BlockParam.basisRound, Maiti2026Power.BlockParam.estN, Maiti2026Power.BlockParam.iHat, Maiti2026Power.BlockParam.thetaHat, Maiti2026Power.normalizeBall, Maiti2026Power.BlockParam.recommend, Maiti2026Power.BlockParam.nextAction, Maiti2026Power.BlockParam.alg, Maiti2026Power.BlockParam.output}\leanok
\uses{def:block_ball_set, def:meta_algorithm, lem:ball_identification, lem:seeded}
Let $\eps \in (0,1]$ and $\delta \in (0,1)$, let $N := N_{\mathrm{norm}}(\eps/4, \delta/(2k))$ be
the budget of the meta-algorithm $\mathrm{NormEst}(\eps/4, \delta/(2k))$
(Definition~\ref{def:meta_algorithm}) and $n_2 := \lceil 4(8d + 48\log(2/\delta))/\eps^2 \rceil$.
The \emph{block algorithm} on $\Xs$ is the seeded algorithm (Lemma~\ref{lem:seeded}) with
uniform sign-vector seeds and budget $N_{\mathrm{sep}}(\eps,\delta) := kN + dn_2$ defined by the
following round schedule.
\begin{enumerate}
  \item[(1)] For $i = 0, \dots, k-1$, the rounds $[iN, (i+1)N)$ run
        $\mathrm{NormEst}(\eps/4, \delta/(2k))$ on block $i$: the action of round $iN + t$ is
        $\iota_i$ applied to the action of round $t$ of the meta-algorithm computed from the
        ($\pi_i$-projected) history of the block and the seed of the round. The estimate
        $\hat r_i$ is the estimate of the meta-algorithm computed from the observations of the
        rounds of block $i$, and $\hat i := \argmax_i\hat r_i$ (smallest index in case of ties).
  \item[(2)] The rounds $kN + mn_2 + \ell$, $m < d$, $\ell < n_2$, play $\iota_{\hat i}(e_m)$
        (the basis design of Lemma~\ref{lem:ball_identification} on block $\hat i$); let
        $\hat\vartheta \in \R^d$ be the vector of the $d$ empirical means of the observations
        at each $e_m$, and recommend $\rec := \iota_{\hat i}\big(\hat\vartheta/\norm{\hat\vartheta}\big)$
        (with $\hat\vartheta/\norm{\hat\vartheta} := 0$ if $\hat\vartheta = 0$).
\end{enumerate}
\end{definition}

\textbf{Lean remark.} The parameters are gathered in a structure \texttt{BlockParam}
($k \ge 1$, $\eps$, $\delta$), with the meta-algorithm parameters
\texttt{BlockParam.P} $= (\eps/4, \delta/(2k))$. The block of a round $t < kN$ is
$\lfloor t/N \rfloor$ (\texttt{blockOf}), the basis vector of a round of phase (2) is given by
\texttt{basisRound}. The estimate of block $m$ (\texttt{estN}) is the estimate of
Definition~\ref{def:meta_algorithm} applied to the observations $y_{mN + r}$; the selected block
$\hat i$ (\texttt{iHat}, defined by \texttt{Nat.find}) and the empirical mean vector
(\texttt{thetaHat}) are measurable functions of the observation sequence, hence so is the
recommendation (\texttt{recommend}, \texttt{output}); the normalization $v \mapsto v/\norm v$
(\texttt{normalizeBall}, with $0 \mapsto 0$) is measurable but not continuous. The seeded
algorithm \texttt{alg} has next-action map \texttt{nextAction}. The recommendation $0$ when
$\hat\vartheta = 0$ replaces the paper's $e_1$; both have regret at most $2\norm{\hat\vartheta - \vartheta}$.

\textbf{Lean remark.} Phase (1) is itself a $k$-fold composition of algorithms with action
type $\Xs$ (each $\iota_i\mathrm{NormEst}$ has action type $\Xs$ by
Lemma~\ref{lem:separation_block_embedding}); the index $\hat i$ is a
measurable function of the phase-(1) history $h$ with values in $\{1,\dots,k\}$, and phase (2)
is the fixed design of Lemma~\ref{lem:ball_identification} transported by $\iota_{\hat i}$,
again with action type $\Xs$: a deterministic algorithm parametrized by $\hat i$, which is
the ``$\mathrm{alg}_2(h)$'' of Lemma~\ref{lem:composition}. The map
$h \mapsto \iota_{\hat i(h)}\mathrm{alg}$ is measurable in $h$ because $\hat i$ takes finitely
many values: on each of the measurable sets $\{\hat i = i\}$ it is the constant algorithm
$\iota_i\mathrm{alg}$, so its policy kernels are finite case distinctions of kernels
(\texttt{Kernel.piecewise}), and Lemma~\ref{lem:composition} applies.

\begin{lemma}[Selection of a good block]\label{lem:block_selection}
\lean{Maiti2026Power.BlockParam.estN_le_iHat, Maiti2026Power.norm_sub_inner_normalizeBall_le, Maiti2026Power.BlockParam.regret_le}\leanok
\uses{def:block_algorithm, lem:block_ball_basic}
Let $\theta \in \R^{kd}$ and let $y$ be an observation sequence such that
$\abs{\hat r_i - \norm{\theta^{(i)}}} \le \eps/4$ for all $i < k$, where $\hat r_i$ is the estimate
of block $i$ computed from $y$. Then, with $\hat i$ the selected block,
\[
  \max_{i<k}\norm{\theta^{(i)}} \le \norm{\theta^{(\hat i)}} + \frac{\eps}{2} ;
\]
and if moreover the empirical mean vector $\hat\vartheta$ computed from $y$ satisfies
$\norm{\hat\vartheta - \theta^{(\hat i)}} \le \eps/4$, then the recommendation
$\rec = \iota_{\hat i}(\hat\vartheta/\norm{\hat\vartheta})$ has simple regret
$r(\rec, \theta) \le \eps$ on $\Xs$.
\end{lemma}

\begin{proof}
\leanok
\uses{def:block_algorithm, lem:block_ball_basic}
With $i_* \in \argmax_i\norm{\theta^{(i)}}$ and $\hat r_{i_*} \le \hat r_{\hat i}$ (definition of
$\hat i$),
$\norm{\theta^{(\hat i)}} \ge \hat r_{\hat i} - \eps/4 \ge \hat r_{i_*} - \eps/4
\ge \norm{\theta^{(i_*)}} - \eps/2$. For the regret, by Lemma~\ref{lem:block_ball_basic} the
support function of $\Xs$ at $\theta$ is $\max_i\norm{\theta^{(i)}}$ and
$\ip{\iota_{\hat i}\rec'}{\theta} = \ip{\rec'}{\theta^{(\hat i)}}$ with
$\rec' := \hat\vartheta/\norm{\hat\vartheta}$; and for every $\vartheta, v \in \R^d$,
$\norm{\vartheta} - \ip{v/\norm v}{\vartheta} \le 2\norm{v - \vartheta}$ (if $v \ne 0$,
$\ip{v/\norm v}{\vartheta} = \norm v - \ip{v/\norm v}{v - \vartheta} \ge \norm v - \norm{v - \vartheta}
\ge \norm\vartheta - 2\norm{v - \vartheta}$; if $v = 0$ the left side is
$\norm\vartheta = \norm{v - \vartheta}$). Hence
$r(\rec,\theta) \le \big(\norm{\theta^{(\hat i)}} + \eps/2\big) - \big(\norm{\theta^{(\hat i)}} - 2\cdot\eps/4\big) = \eps$.
\end{proof}

\begin{proof}
\uses{thm:norm_estimation, lem:separation_block_embedding, lem:composition}
By Lemma~\ref{lem:separation_block_embedding} and Theorem~\ref{thm:norm_estimation}, the
run on block $i$ is a run of $\mathrm{NormEst}(\eps/4,\delta/(2k))$ against $\nu_{\theta^{(i)}}$ on
$\ball_d$, so $\Prob(|\hat r_i - \norm{\theta^{(i)}}| > \eps/4) \le \delta/(2k)$; by
Lemma~\ref{lem:composition} this bound holds conditionally on the history of the earlier
blocks, hence unconditionally. A union bound over the $k$ blocks gives the first claim with
probability at least $1 - \delta/2$. On this event, with $i_* \in \argmax_i\norm{\theta^{(i)}}$,
\[
  \norm{\theta^{(\hat i)}} \ge \hat r_{\hat i} - \frac\eps4 \ge \hat r_{i_*} - \frac\eps4
  \ge \norm{\theta^{(i_*)}} - \frac\eps2 .
\]
\end{proof}

\begin{theorem}[Polynomial separation; Theorem~7 of the paper]\label{thm:separation}
\lean{Maiti2026Power.BlockParam.seedAction_of_T₂_le, Maiti2026Power.BlockParam.thetaHat_eq, Maiti2026Power.BlockParam.recommend_yOf_seedFinHist, Maiti2026Power.BlockParam.one_sub_le_seedMeasure_real_output, Maiti2026Power.BlockParam.T_le}\leanok
\uses{def:block_algorithm, def:pac, thm:separation_lower_nonadaptive, lem:seeded}
Let $\eps \in (0,1]$ and $\delta \in (0,1)$. The block algorithm is $(\eps,\delta)$-PAC on $\Xs$
(with failure probability at most $15\delta/16$), with budget
\[
  N_{\mathrm{sep}}(\eps,\delta)
  \le \frac{1600192\,kd\log(8k/\delta) + 33\,d^2}{\eps^2}
  \le C_{\mathrm{sep}}\,\frac{kd\log(8k/\delta) + d^2}{\eps^2},
  \qquad C_{\mathrm{sep}} := 2\cdot 10^6 .
\]
If $d \le k$, this is at most $2C_{\mathrm{sep}}\,kd\big(\log(1/\delta) + \log(8k)\big)/\eps^2$,
whereas by Theorem~\ref{thm:separation_lower_nonadaptive} every non-adaptive
$(\eps,\delta)$-PAC algorithm with $\delta < 1/16$ needs at least
$\tfrac12\big(c_{\mathrm{ad}}\,kd\log(1/\delta) + kd^2/210\big)/\eps^2$ samples. Here
``non-adaptive'' means a deterministic fixed design with an arbitrary (possibly randomized)
recommendation kernel, the class of Definition~\ref{def:fixed_design} and of
Theorem~\ref{thm:lower_nonadaptive}; see Remark~\ref{rem:separation_gap}.
\end{theorem}

\begin{corollary}[Theorem~7 of the paper, existential form of the upper bound]\label{cor:separation_exists}
\lean{Maiti2026Power.exists_isPAC_blockBallSet}\leanok
\uses{thm:separation, def:pac}
Let $k \ge 1$, $\eps \in (0,1]$ and $\delta \in (0,1)$. There exist a budget
$T \le C_{\mathrm{sep}}\,(kd\log(8k/\delta) + d^2)/\eps^2$ and an identification algorithm with
budget $T$ that is $(\eps,\delta)$-PAC on the block-ball set $\Xs \subset \R^{kd}$.
\end{corollary}
\begin{proof}
\leanok
\uses{thm:separation, lem:seeded}
Take the block algorithm of Theorem~\ref{thm:separation}. (For $d = 0$ every reward vector is
$0$, every recommendation is optimal, and the budget $0$ suffices.)
\end{proof}

\begin{proof}
\leanok
\uses{lem:block_selection, lem:ball_identification, lem:separation_block_embedding, lem:seeded, lem:block_ball_basic, thm:norm_estimation, thm:separation_lower_nonadaptive, lem:statistic_structure}
\emph{Correctness.} By Lemma~\ref{lem:seeded} it suffices to prove the bound on the seed space
$(\Omega_{\mathrm{seed}}, \Prob_{\mathrm{seed}})$ of Lemma~\ref{lem:statistic_structure}, where
the recommendation of the run is the recommendation computed from the observation sequence
(it only reads the rounds $t < N_{\mathrm{sep}}$). Consider the two failure events
\[
  E_1 := \bigcup_{i<k}\Big\{ \abs{\hat r_i - \norm{\theta^{(i)}}} > \frac\eps4 \Big\},
  \qquad
  E_2 := \Big\{ \norm{\Delta} > \frac\eps4 \Big\},
  \quad \Delta := \pi'_{kN, n_2}\text{-average of the noises},
\]
where $\Delta \in \R^d$ has coordinates $\Delta_m := \frac1{n_2}\sum_{\ell<n_2}\eta_{kN + mn_2 + \ell}$.
By Lemma~\ref{lem:separation_block_embedding}, $\hat r_i$ is the estimate of
$\mathrm{NormEst}(\eps/4, \delta/(2k))$ against $\nu_{\theta^{(i)}}$ on the shifted seed sequence,
whose law is $\Prob_{\mathrm{seed}}$, so by Theorem~\ref{thm:norm_estimation} each event of the
union has probability at most $\tfrac78\cdot\delta/(2k)$ and $\Prob_{\mathrm{seed}}(E_1) \le 7\delta/16$.
In phase (2) the actions are $\iota_{\hat i}(e_m)$, so the empirical mean vector is
$\hat\vartheta = \theta^{(\hat i)} + \Delta$, where $\Delta$ does not depend on $\hat i$; the law of
the noises of the basis window is $\Normal(0,1)^{\otimes dn_2}$
(Lemma~\ref{lem:statistic_structure}), so by the chi-square tail bound of
Lemma~\ref{lem:ball_identification} with $c = \eps/4$ and the choice of $n_2$
($n_2\eps^2/16 \ge 2d + 12\log(2/\delta)$), $\Prob_{\mathrm{seed}}(E_2) \le \delta/2$. Outside
$E_1 \cup E_2$, Lemma~\ref{lem:block_selection} gives $r(\rec,\theta) \le \eps$. Hence
$\Prob_{\mathrm{seed}}(r(\rec,\theta) \le \eps) \ge 1 - 7\delta/16 - \delta/2 = 1 - 15\delta/16 \ge 1 - \delta$.
No conditioning on the first phase is needed: the second failure event only involves the
noises of the basis window, whatever the selected block.

\emph{Budget.} By Theorem~\ref{thm:norm_estimation} with $(\eps/4, \delta/(2k))$
(note $\log(4/(\delta/(2k))) = \log(8k/\delta)$), each run of
phase (1) uses at most $10^5\,d\log(8k/\delta)\cdot 16/\eps^2$ rounds, so phase (1) uses at
most $1.6\cdot10^6\,kd\log(8k/\delta)/\eps^2$ rounds. Phase (2) uses
$dn_2 \le d\big( (32d + 192\log(2/\delta))/\eps^2 + 1 \big) \le (33d^2 + 192\,d\log(2/\delta))/\eps^2$
(using $d \le d^2/\eps^2$), and $192\,d\log(2/\delta) \le 192\,kd\log(8k/\delta)$. Summing gives
the first bound, and $1600192 \le C_{\mathrm{sep}}$, $33 \le C_{\mathrm{sep}}$ give the second. If
$d \le k$ then $d^2 \le kd \le kd\log(8k/\delta)$ (as $\log(8k/\delta) \ge \log 8 > 1$), whence
the third bound with $\log(8k/\delta) = \log(1/\delta) + \log(8k)$. The lower bound is
Theorem~\ref{thm:separation_lower_nonadaptive}.
\end{proof}

\begin{remark}[The size of the gap]\label{rem:separation_gap}
For $d \le k \le d^2$ and a fixed $\delta$, the adaptive budget is
$O(kd\log k/\eps^2) = O(kd\log d/\eps^2)$ (since $\log k \le 2\log d$) while every
non-adaptive algorithm needs $\Omega(kd^2/\eps^2)$: adaptivity gains a factor
$\Omega(d/\log d)$, polynomial in the ambient dimension $kd \le d^3$. The paper phrases the
lower bound as $\Omega(kd\log(1/\delta)/\eps^2 + kd^2/\eps^2)$ using the adaptive lower bound
of Chen et al.\ \cite{chen2024assouad} on each block together with a padding argument for
expected budgets; the route above through Theorem~\ref{thm:lower_nonadaptive} gives the same
$kd^2$ term with the explicit constant $1/210$ for deterministic budgets and $\delta \le 1/10$,
without any citation.

\emph{Restriction of the comparison class.} The paper's Theorem~7 states the lower bound for
``possibly randomized non-adaptive'' algorithms. The lower bound proved here
(Theorem~\ref{thm:separation_lower_nonadaptive}(i)) covers only \emph{deterministic} fixed
designs, with an arbitrary randomized recommendation rule: this is the class of
Definition~\ref{def:fixed_design} and of Theorem~\ref{thm:lower_nonadaptive}, and it is the
class in which Theorem~\ref{thm:separation} compares the block algorithm with non-adaptive
ones. A non-adaptive algorithm whose design is drawn at random (independently of the
observations) is not covered: for such an algorithm the pigeonhole step selects a
block depending on the realized design, and the fixed-design lower bound would have to be
applied conditionally on the design and integrated, which is not done here. The adaptive
lower bound (ii) applies to every algorithm, randomized designs included, but only gives the
$kd\log(1/\delta)$ term. This restriction is recorded in the list of deviations from the
paper in the overview.

The constant $C_{\mathrm{sep}}$ inherits the (unoptimized) constant
$C_{\mathrm{norm}}$ of Theorem~\ref{thm:norm_estimation} multiplied by $16$ because the
norm-estimation accuracy is $\eps/4$.
\end{remark}
